% Part: normal-modal-logic
% Chapter: tableaux
% Section: introduction

\documentclass[../../../include/open-logic-section]{subfiles}

\begin{document}

\olfileid[hi]{nml}{tab}{int}

\olsection{परिचय}

!!^{tableau}, !!{signed formula}ों के कुछ (नीचे की ओर शाखित) वृक्ष हैं;
अर्थात् सत्यता-मूल्य चिह्न ($\True$ या $\False$) और किसी !!a{sentence} से
बने युग्म
\[
\sFmla{\True}{!A} \text{ या } \sFmla{\False}{!A}.
\]
!!^a{tableau} कुछ \emph{मान्यताओं} से आरंभ होता है। प्रत्येक अगला
!!{signed formula} किसी अनुमान-नियम को लगाकर उत्पन्न होता है। कुछ
अनुमान-नियम वृक्ष के किसी सिरे पर एक या अधिक !!{signed formula} जोड़ते हैं;
अन्य दो नए सिरे जोड़ते हैं, जिससे दो शाखाएँ बनती हैं। नियम ऐसे !!{signed formula} देते हैं जिनका !!{formula} उस !!{signed formula} के सूत्र से कम जटिल
होता है जिस पर नियम लगाया गया था। जब किसी शाखा में
$\sFmla{\True}{!A}$ और $\sFmla{\False}{!A}$, दोनों हों, तो हम उस शाखा को
\emph{संवृत} कहते हैं। यदि किसी !!a{tableau} की हर शाखा संवृत हो, तो पूरा
!!{tableau} संवृत है। कोई संवृत !!{tableau} ऐसी !!a{derivation} बनाता है जो
दिखाती है कि टैब्लो आरंभ करने के लिए प्रयुक्त !!{signed formula}ों का
समुच्चय असंतोषणीय है। इससे $\Proves$ संबंध परिभाषित किया जा सकता है:
$\Gamma \Proves !A$ तभी जब कोई परिमित समुच्चय~$\Gamma_0 = \{!B_1,
\dots, !B_n\} \subseteq \Gamma$ हो ऐसा कि इन मान्यताओं के लिए कोई संवृत
!!{tableau} हो:
\[
\{\sFmla{\False}{!A}, \sFmla{\True}{!B_1}, \dots, \sFmla{\True}{!B_n}\}.
\]

मोडल तर्कशास्त्रों के लिए हमें !!{signed formula} की धारणा का विस्तार भी
करना है और ऐसे नियम भी जोड़ने हैं जो
\iftag{prvBox}{$\Box$\iftag{prvDiamond}{ और
    $\Diamond$}}{$\Diamond$} को समेटें। किसी चिह्न ($\True$ या $\False$)
के अतिरिक्त, मोडल !!{tableau} के !!{formula}ों में \emph{उपसर्ग}~$\sigma$
भी होते हैं। उपसर्ग धन पूर्णांकों के अशून्य अनुक्रम हैं, अर्थात् $\sigma
\in (\PosInt)^* \setminus \{\emptyseq\}$। हम ऐसे उपसर्गों को आसपास के
$\tuple{\ }$ के बिना लिखते हैं और अलग-अलग !!{element}ों को अल्पविराम के
स्थान पर बिंदुओं से अलग करते हैं। यदि $\sigma$ कोई उपसर्ग है, तो
$\sigma.n$, $\sigma \concat \tuple{n}$ है; जैसे यदि $\sigma = 1.2.1$,
तो $\sigma.3$, $1.2.1.3$ है। अतः उदाहरणार्थ,
\[
\sFmla{\True}{\Box !A \lif !A}[1.2]
\]
एक \emph{उपसर्गयुक्त !!{signed formula}} है (या संक्षेप में केवल
\emph{उपसर्गयुक्त !!{formula}})।

सहजतः, उपसर्ग किसी प्रतिरूप में ऐसे जगत को नाम देता है जो किसी !!a{tableau}
की शाखा के !!{formula}ों को संतुष्ट कर सकता है; और यदि $\sigma$ किसी जगत को
नाम देता है, तो $\sigma.n$, $\sigma$ द्वारा नामित जगत से अभिगम्य किसी
जगत को नाम देता है।

\end{document}
